\subsection{Program Analysis for Smart Contracts}

% \noindent \textbf{Parametric Optimization.}
% For some flash-loan attack cases, prior work manually extracted formulas of candidate functions,
% manually defined parameter constraints, and applied off-the-shelf optimization to search for profitable parameters\cite{cao2021flashot,qin2021attacking}.
% These methods demonstrate feasibility for specific incidents, but their manual workflow and protocol-specific expertise requirements
% limit scalability when the attack space is large.

\noindent \textbf{Static Analysis and Symbolic Execution.}
Static analyzers such as Slither, Securify, Zeus, Park, and 
SmartCheck\cite{feist2019slither,tsankov2018securify,
kalra2018zeus,zheng2022park,tikhomirov2018smartcheck}
have been widely used to detect vulnerabilities in smart contracts.
Symbolic-execution-based tools and 
techniques\cite{king1976symbolic,mythril,luu2016making,liu2020towards,
frank2020ethbmc,feng2019precise,mossberg2019manticore}
systematically explore path conditions to find property-violating traces.
These approaches are highly effective for many vulnerability classes. 
However, for complex DeFi smart contracts, real-world attacks typically 
involves multiple contracts and long execution traces, 
which can lead to state coupling and path explosion 
challenges for static and symbolic methods.

\noindent \textbf{Fuzzing.}
ContractFuzzer, sFuzz, ContraMaster, SMARTIAN, 
and ItyFuzz\cite{jiang2018contractfuzzer,nguyen2020sfuzz,wang2020oracle,
choi2021smartian,shou2023ityfuzz}
introduce effective fuzzing strategies for smart-contract security testing.
However, these techniques either only work on one contract or focus
on specific vulnerabilities like re-entrancy. 
Moreover, real-world DeFi smart contract attacks can 
involve tens of contract interactions and tens of chosen parameters,
which is unlikely to be found by random fuzzing. 


\noindent \textbf{Smart Contract Attack Monitoring.}
Prior studies analyze DeFi smart contract attacks and related
 DeFi transaction strategies\cite{qin2021attacking,cao2021flashot,
 werapun2022flash,wang2021towards,gan2022understanding}.
Monitoring-oriented systems have also been proposed 
for post-hoc or real-time attack detection
\cite{xia2023detecting,ramezany2023midnight,wang2022defiscanner}.
Despite these efforts and the good coverage of known attack patterns, 
these works do not propose systematic generalized methods for defending 
against novel attack vectors that may arise in the future. 

